\documentclass[10pt]{article}

\usepackage[utf8]{inputenc}

\usepackage[T1]{fontenc}

\usepackage{amsmath}

\usepackage{amsfonts}

\usepackage{amssymb}

\usepackage[version=4]{mhchem}

\usepackage{stmaryrd}

\usepackage{bbold}

\usepackage{graphicx}

\usepackage[export]{adjustbox}

\graphicspath{ {./images/} }



\begin{document}

ненулевой $p$-вектор $x_{1} \wedge \ldots \wedge x_{p} ;$ такие П. ваз. р а вло жи м Ы м или п р о с т Ы м и (часто-просто П.). При этом линейно независимые системы $x_{1}$, .., $x_{p}$ и $y_{1}, \ldots ., y_{p}$ порождают одно и то же подпространство в $V$ в том и только в том случае, когда $y_{1} \wedge$ $\cdots \wedge y_{p}=c x_{1} \wedge \cdots \wedge x_{p}$, где $c \in K$. Для любого ненулевого поливектора $t \in \Lambda^{p} V$ его а н н у л я т о р Ann $t=\{v \in V \mid t \wedge v=0\}$ есть подпространство размерности $\leqslant p$, причем поливектор $t$ разложим тогда и только тогда, когда $\operatorname{dim} \operatorname{Ann} t=p$. Разложимые $p$ векторы $n$-мерного пространства $V$ образуіт коническое алгебраич. многообразие в $\Lambda^{p} V$; соответствующее проективное алгебраич. многообразие есть Грассмана многообразие. Любой ненулевой $n$-вектор или $(n-1)$ вектор в $n$-мерном пространстве $V$ разложим; бивектор $t$ разложим тогда и только тогда, когда $t \wedge t=0$. Если $v_{1}, \ldots, v_{n}$ - базис пространства $V$ и $x_{i}=$ $=\sum_{j=1}^{n} x_{i}^{j_{j}} v_{j}$, то координатами полівектора $t=x_{1} \wedge \ldots$ $\wedge x_{p}$ в базисе $\left\{v_{i_{1}} \wedge \ldots \wedge v_{i_{p}} \mid i_{1}<\ldots<i_{p}\right\}$ пространства $\Lambda^{p} V \quad$ являются миноры $t^{i_{1} \ldots i^{\prime}} p=\operatorname{det}\left\|x_{i}^{j_{k}}\right\|$, $i_{1}<\ldots<i_{p}$, матрицы $\left\|x_{i}^{j}\right\|$. В частности, при $p=n$



\[

x_{1} \wedge \ldots \wedge x_{n}=\operatorname{det}\left\|x_{i}^{j}\right\| v_{1} \wedge \ldots \wedge v_{n} .

\]



Если фиксировать ненулевой $n$-вектор $\omega \in \Lambda^{n} V$, то возникает двойственность между р-векторами и $(n-p)$-векторами, т. е. естественный изоморфизм



\[

\pi: \Lambda^{p}(V) \longrightarrow\left(\Lambda^{n-p} V\right)^{*} \cong \Lambda^{n-p}\left(V^{*}\right)

\]



такой, что $t \wedge u=\pi(t)(u) \omega$ для всех $t \in \Lambda^{p} V$ и $u \in \Lambda^{n-p} V$.



Пусть $k=\mathbb{R}$ и в $V$ задано скалярное произведение, тогда в $\Lambda^{p} V$ индуцируется скалярное произведение, обладающее следующим свойством: для любого ортонормированного базиса $v_{1}, \ldots, v_{n}$ в $V$ базис $\left\{v_{i_{1}} \wedge \ldots\right.$. $\left.\wedge v_{i_{p}} \mid i_{1}<\ldots<i_{p}\right\}$ в $\Lambda^{p} V$ также ортонормирован. Скаля рный квадрат



\[

(t, t)=\sum_{i_{1}<\ldots<i_{p}}\left(t^{i_{1} \ldots i_{p}}\right)^{2}

\]



разложимого поливектора $t=x_{1} \wedge \ldots \wedge x_{p}$ совпадает с квадратом объема параллелепипеда в $V$; натянутого на векторы $x_{1}, \ldots, x_{p}$. Если фиксировать в $n$-мерном евклидовом пространстве $V$ ориентацию (что равносильно выбору $n$-вектора $\omega$, для к-рого $(\omega, \omega)=1)$, то указанная выше двойственность приводит к естественному изоморфизму $\gamma: \Lambda^{p} V \rightarrow \Lambda^{n-p} V$ такому, что $t \wedge u=$ $=(\gamma(t), u) \omega$ для всех $t \in \Lambda^{p} V, u \in \Lambda^{n-p} V$. В частности, $(n-1)$-вектору $t=x_{1} \wedge \ldots \wedge x_{n-1}$ соответствует вектор $\gamma(t) \in V$, наз. в е к т о р ным п роиз веден и е м векторов $x_{1}, \ldots, x_{n-1}$.



Лum.: [1] Бур баки Н., Алгебра. Алгебраические структуры. Линейная и полилинейная алгебра, пер. с франц. М., 1962; [2] К о с т р и к и н А. И., М а н и н Ю. И., Јинейная алгебра и геометрия, М., 1980; [3] П о с т н и к о в М. М., Линейная алгебра и дифференциальная геометрия, М., 1979.



ПОЛИГАРМОНИЧЕСКАЯ ФУНКЦИЯ, Г п п рг а моническая ф ф у ниция, ме та га р-



\includegraphics[max width=\textwidth, center]{2023_07_08_68d09d27314fe394f4a1g-1}

ция $u(x)=u\left(x_{1}\right.$, . .., $\left.x_{n}\right)$ действительных переменных, огределенная в области $D$ евклидова пространства $\mathbb{R}^{n}, \quad n \geqslant 2$, имеющая непрерывные частные производные до 2 -го порядка включительно и удовлетворяющая всюду в $D$ по ли г а р м о ни ч е с к м у у р а в в е и ю



\[

\Delta^{m} u \equiv \Delta(\Delta \ldots(\Delta u))=0, \quad m \geqslant 1,

\]



где $\Delta$-оператор Лапласа. При $m=1$ получаются гармонические бункции, при $m=2$ - бигармониче- ские функции. Каждая П. Ф. есть аналитич. Функция от координат $x_{j}$. Нек-рые другие свойства lармонич. функций также переносятся с соответствуюшими изменениями на П. ф.



Для П. ф. любого порядка $m>1$ обобцаются представления при помощи гармонич. функций, известные для бигармонич. функций (см. [1] - [5]). Напр., для П. Ф. $u$ двух переменных сшраведливо представление



\[

u\left(x_{1}, x_{2}\right)=\sum_{k=0}^{m-1} r^{2 k} \omega_{k}\left(x_{1}, x_{2}\right), r^{2}=x_{1}^{2}+x_{2}^{2},

\]



где $\omega_{k}, k=0, \ldots, m-1$ - гармонич. функции в области $D$. Для того чтобы функция $u\left(x_{1}, x_{2}\right)$ двух переменных была П. ф., необходимо и достаточно, чтобы она была действительной (или мнимой) частью полианалитической функции.



Основная краевая задача для П. ф. порядка $m>1$ состоит в следующем: найти П. ф. $u=u(x)$ в области $D$, непрерывную вместе с производными до $(m-1)$-го порядка включительно в замкнутой области $\vec{D}=$ $\equiv D \cup C$ и удовлетворяющую на границе $C$ условиям:



\[

\begin{aligned}

\left.u\right|_{C} & =f_{0}(y), \\

\left.\frac{\partial u}{\partial n}\right|_{C} & \left.=f_{1}(y), \ldots,\left.\frac{\partial^{m-1} u}{\partial n^{m-1}}\right|_{C}=m-1(y), y \in C, \quad\right\}

\end{aligned}

\]



где $\partial u / \partial n-$ производная по нормали к $C$, а $f_{0}(y)$, ..., $f_{m-1}(y)$ - заданные достаточно гладкие функции на достаточно гладкой границе $C$. Многие исследования посвяшены решению задачи (\textit{) в шаре пространства $\mathbb{R}^{n}$ (см. [1], [6]). Для решения задачи (}) в случае произвольной области применялся метод интегральных уравнений, а также различные вариационные методы (см. [1], [6]).



Лит.: [1] В е к у а И. Н., Новые методы решения эллиптических уравнений, М.-Ј., 1948; [2] П р и в а л о в И. И., [3] $\mathrm{N}$ i c o l e s c o $\mathrm{M}$. Les fonctions polyharmoniques, $\mathrm{P}-1936$ [4] e r o ж e, "Disq. Math. Phys.", 1940 , v. 1, p. 43-56; [5] T. 0 lo t t i C., "Giorn. Math. Battaglini", 1947, v. 1, p. 61-117 [6] М и р а н д а К., Уравнения с частными производными ПОЛИГОН над моноидом $R, R-\square$ л иг о н, о п еp а в д,- непустое множество с моноидом операторов. Точнее, непустое множество $A$ наз. ле в ы м П. над моноидом $R$, если для любых $\lambda \in R$ и $a \in A$ определено произведение $\lambda a \in A$, причем



$\mathbf{\Lambda}$



\[

\begin{aligned}

(\lambda \mu) a & =\lambda(\mu a) \\

1 a & =a

\end{aligned}

\]



для любых $\lambda, \mu \in R, a \in A . \Pi$ р а в ы й П. определяется аналогично. Задание $R$-полигона $A$ равносильно заданию гомоморфизма $\varphi$ моноида $R$ в моноид отображений множества $A$ в себя, переводящего 1 в тождественное птображкение. При этом $\lambda a=b$ тогда и только тогда, когда



\[

\varphi(\lambda)(a)=b .

\]



В qастности, каждое непустое множество можно рассматривать как П. над моноидом его отображений в себя. Таким образом, П. тесно связан с представлением полугруппы преобразованиями.



Если $A$ - универсальная алгебра, сигнатура к-рой $\Omega$ содерякит лишь унарные операции, то $A$ можно превратить в П. над свободным моноидом $F$ с системой свободных образующих $\Omega$, положив



\[

\left(f_{1} f_{2} \ldots f_{n}\right) a=f_{1}\left(f_{2}\left(\ldots\left(f_{n}(a)\right) \ldots\right)\right)

\]



для любых $f_{i} \in \Omega, a \in A$. Если $\Omega-$ множество входных сигналов автомата с множеством состояний $A$, то $A$ аналогичным образом превращается в $F$-полигон (cp. Автоматов алгебраическая теория).



Отображение $\varphi R$-полигона $A$ в $R$-полигон $B$ наз. гомоморфи змом, если $\varphi(\lambda a)=\lambda \varphi(a)$ для лю-





\end{document}